\section{A resolved information and inference design for servicing}
\label{s6:servicing}
The privately logged three-state model is useful precisely because it can
fail despite an identified population value. Resolving that failure requires
separating the statistical experiment from the optimizer and from the
proposed decision. Here the servicing decision is whether to offer temporary
assistance; the objective is discounted collected cash net of assistance and
operating cost. Signals arrive before the decision, while the cash rewards
used in this four-decision policy are settled after its decision horizon.
The policy may use signals and its own previous actions. It may not use the
later audited state or settled cash as contemporaneous information.

\subsection{What the retained failure actually shows}
For the original 600-trajectory experiment, all nine unavailable fits were
on a constrained transition-logit boundary. Reconstructing those exact seeds
reproduces the stored values, rather than diagnosing a different simulation.
With 90 emission labels, 76 of 80 fits were available: conditional coverage
was .8947, empirical value standard deviation 10.668 and mean model standard
error 9.657. With 600 labels, 75 were available: conditional coverage was
.9333, empirical standard deviation 9.141 and mean standard error 9.687.
Counting unavailable fits gives the previously reported .850 and .875.
The latter case therefore does not mainly show underestimated standard
errors among the usable fits. Availability, tail error and bias are different
questions. Eighty repetitions leave substantial Monte Carlo uncertainty.

The binding coordinates are transition logits, including near-zero
destinations under privately selected actions. Larger emission calibration
alone does not supply those scarce action/state transitions. The model's
positive population transition probabilities do not prevent an observed
likelihood maximum near the boundary. The observed fit reconstructions,
larger-trajectory follow-up and complete-state frequency experiment separate
this scarcity explanation from the earlier differencing/start checks. They
do not certify a global maximum or establish one universal cause of
undercoverage. At 3,000 observational trajectories and 600 independent
emission labels, all 120 follow-up fits are available; coverage is .9417,
bias $-.213$, empirical standard deviation 4.097 and mean reported standard
error 4.301. This supports a bounded larger-sample regime on this mechanism,
not a universal 3,000-customer rule. The failed 199-refit bootstrap remains
adverse evidence for the original design.

\subsection{A measurement acquisition alternative}
An independent representative audit now follows whole customer trajectories,
recording the state before each action, the same production signal, the
actual action and the eventually settled reward. The states must have an
operational definition: for example, a predecision classification of verified
repayment capacity and available liquidity into solvent, temporarily
constrained and persistent-shortfall categories. Definition and thresholds
are fixed before outcomes. Independent adjudication checks timestamped
income, obligations and accessible balances; disagreements and unavailable
documents remain measurement failures.

These categories are not automatically sufficient causal states. The model
asserts that, conditional on the state, the officer's private action
information adds no further outcome/transition cause; that the state is
Markov for this horizon; and that the audit and production emissions agree.
A coarse affordability label can violate all three. Audit selection must
represent the target population, or known inclusion probabilities and their
influence must enter the estimator. The implementation below assumes
representative independent trajectories. It does not silently replace a
convenience audit by a random sample.

Let $C_{i,sj}^E$ count signal $j$ in state $s$ along audited path $i$,
$C_{i,asj}^T$ count transitions from $s$ to $j$ under action $a$, and
$N_{i,as},R_{i,as}$ be state/action exposure and reward sums. The row-frequency
and reward estimators are
\[
 \hat E_{sj}=\frac{\sum_i C_{i,sj}^E}{\sum_{i,j'}C_{i,sj'}^E},\quad
 \hat T_{asj}=\frac{\sum_i C_{i,asj}^T}{\sum_{i,j'}C_{i,asj'}^T},\quad
 \hat r_{as}=\frac{\sum_iR_{i,as}}{\sum_iN_{i,as}}.
\]
Initial frequencies estimate $\mu$; action frequencies estimate the private
logger $q$. This is the complete-data likelihood estimate in the declared
finite model. It has no local mixture-likelihood optimization. It buys that
simplicity with a much more expensive observation design.

For any row with path counts $C_{ij}$, let $c_j=\E C_{ij}>0$ and use
logit coordinates $\ell_j=\log(c_j/c_K)$. Direct differentiation gives
\[
 \varphi_{ij}^{\ell}=C_{ij}/c_j-C_{iK}/c_K,\qquad
 \varphi_{i,as}^{r}=(R_{i,as}-r_{as}N_{i,as})/\E N_{i,as}.
\]
Stack these with the initial and emission coordinates. Their path-level
outer products preserve all within-trajectory covariance; treating visits
as independent would not. For the scaled reward coordinates divide the
reward influence by the known reward standard deviation. With fixed horizon,
positive row masses, finite second moments and accurate representative
labels, the multivariate iid central limit theorem and quotient delta method
give
\[
 \sqrt n(\hat\vartheta-\vartheta_0)
 =n^{-1/2}\sum_i\varphi_i+o_p(1),\qquad
 \widehat{\Var}(\hat V)=
 \nabla V(\hat\vartheta)'\left(n^{-2}\sum_i\hat\varphi_i\hat\varphi_i'\right)
 \nabla V(\hat\vartheta).
\]
An empty category makes these interior coordinates unavailable. Arbitrary
pseudocounts would define a different finite procedure; they are not added
to disguise support failure. This is a pointwise regular inference route,
not uniform coverage near vanishing state/action probabilities.

\subsection{A dynamic observable policy}
For a prior belief $b$ immediately before signal $o$, update
$b_o(s)=b(s)E_{so}/\sum_jb(j)E_{jo}$. If action $a$ is intervened on, the
next prior is $b_oT_a$. There is no private logger factor in that update:
the policy chose the action. For remaining horizon $h$,
\begin{equation}
 V_h(b)=\sum_oP(o\mid b)\max_a
       \{b_or_a+\gamma V_{h-1}(b_oT_a)\},\qquad V_0=0.
 \label{s6:belief-policy}
\end{equation}
The finite recursion is the belief-state planning construction
\citep{kaelbling1998}. Four signals, two actions and four decisions permit
exact tree enumeration here. The tree is optimized using an independent
2,000-trajectory audit pilot, then frozen in its entirety. Evaluation fits
kernels on separate audited trajectories and propagates their uncertainty
through this fixed tree. It differentiates neither the argmax nor the pilot
training procedure. The interval targets the value of the policy that was
actually trained, conditional on that pilot, not the unknown optimal value.

To evaluate the fixed tree without numerical posterior divisions, propagate
unnormalized state mass. At each signal branch multiply by $E_{\cdot o}$,
add the selected action's expected reward, and propagate through $T_a$.
Summing all leaves gives a polynomial in the kernels and initial law, linear
in the rewards. Its derivative is therefore well-defined for the fixed tree;
central differencing is checked against a smaller step and independent
likelihood/count perturbations. Exact enumeration avoids a filter
approximation error in this selected problem. Larger horizons require
pruning or approximate planning and an explicit value-loss bound.

On the common mechanism the frozen learned policy has true four-period
contrast 25.029 against never assisting. Some actions depend on earlier
signals even after holding the current signal and remaining horizon fixed;
this is a belief-history policy, rather than an immediate-reward rule.
At 1,200 audited evaluation trajectories, all 400 fits are available and
coverage is .950 for this contrast. The original eight-period fixed-signal
contrast has coverage .9425. At 300 audited paths, 54 of 400 fits have an
empty category; all-request coverage is .8275 for the learned policy even
though coverage conditional on availability is .9566. Measurement-supported
identification therefore still requires enough informative visits.

\subsection{Three information routes on the same servicing problem}
An observed-history g-formula substitutes the factual conditional law given
public signals and actions. Under private logging its state posterior after
action $a$ is proportional to $b_o(s)q_s(a)$, while the intervention posterior
is $b_o(s)$. Thus even an exactly estimated public-history conditional law
uses the wrong state mixture for the intervention. Adding a flexible learner
cannot remove this information gap. Recording the action-relevant private
information, or randomizing assistance with adequate support, changes that
conclusion. The common-mechanism comparison evaluates both recursions on the
same frozen policy. Its exact public-history g-formula contrast is $-35.647$,
against the intervention contrast $25.029$. The sign reversal occurs even
without nuisance estimation error.

Validated emissions can identify the kernels using the earlier marked
moment/hidden-likelihood route, with fewer labels but difficult finite
curvature and boundary behavior. Fully audited trajectories identify the
same kernels and policy target more directly, with the explicit acquisition
cost above. Their precision and availability are properties of different
sampling designs, not proof that hidden-state modeling is unnecessary.

Proximal evaluation instead seeks policy value without recovering named
states. For this servicing mechanism an action-side proxy might be a
separately recorded officer assessment, and an outcome-side proxy an
independent affordability measurement. Merely calling them proxies proves
nothing: after conditioning on the latent confounder and allowed context,
the action-side proxy and logged action must be independent of the
outcome-side proxy, proposed action and counterfactual continuation under
the source's hybrid laws. The proposed policy cannot reuse forbidden private
proxy information. \citet[Assumptions 1--2, Theorem 2 and Table 1]{bennettprox2023}
give these policy-dependent restrictions and the recursive action/outcome
bridges. In particular their outcome bridge targets
$\E[\mathbf1\{E_t=A_t\}Y_t\mid U_t,A_t=a]$, where $E_t$ denotes the
proposed action; dropping that indicator changes the theorem.

The source's recursive bridge score is a substantive challenger when the
bank can defend those exclusions and bridge existence/rank conditions but
cannot obtain accurate state labels. The current production feed has one
public signal stream and no independently validated second proxy. It cannot
invoke that theorem by relabelling the same signal twice. If an additional
proxy is collected, its conditional operator, range, regularization and
policy information set become new scientific inputs. This explains why the
audited trajectory route is selected for the bounded decision capability
and proximal evaluation remains a qualified alternative, rather than
claiming that one existing dataset satisfies all three information regimes.

\subsection{The initial product boundary}
The offering includes finite-state estimation and observable-history policy
evaluation under a qualified measurement design, plus the earlier exact
linear-Gaussian controlled likelihood. It does not include arbitrary nonlinear
private-state fitting with routinely calibrated policy intervals. Choosing
an extended Kalman or particle filter would require a declared likelihood,
initial/noise laws, derivative accuracy and policy-value comparison against
a numerical reference, following the approximation obligations already
derived. Section~\ref{c7:nonlinear} supplies a selected audited nonlinear
comparison with additional information and policy-class restrictions.
Those restrictions affect which bank requests can be accepted; they are
not postponed as a generic deployment check. A noisy or selected audit
also changes the statistical model, even when all matrices are well-conditioned.
In a separately registered 200-repetition sensitivity run, independently
misclassifying 10\% of predecision state labels at 1,200 audit trajectories
produces bias $-105.72$ for the original $59.56$ contrast and zero coverage.
This is a specific severe misspecification example, not a calibrated bound
for every 10\% audit-error process.
